% ============================================================================
% MARKOVTRIE: ADAPTIVE PROBABILISTIC LEARNING WITHOUT GRADIENTS
% ============================================================================
\subsection{The MarkovTrie Store}
\label{sec:markovtrie}

The MarkovTrie is the system's learning substrate: an adaptive
probabilistic suffix trie that learns online from every observation
without training epochs, loss functions, or weight matrices.  Each node
in the Kadabra routing layer (\Cref{sec:kadabra_dht}) owns one
MarkovTrie instance; Values are routed to tries by affinity similarity,
so each trie naturally specializes in a region of content space.

\subsubsection{Data Structure}

The trie is a rooted tree in which each node $n$ stores:
\begin{itemize}[nosep]
  \item A token $\tau_n$ (the symbol at this position in the suffix).
  \item A children map $\mathcal{C}_n : \mathcal{V} \to \mathcal{N}$
        from vocabulary symbols to child nodes.
  \item Per-label decayed counts
        $c_n(\ell) \in \mathbb{R}_{\geq 0}$ for each label
        $\ell \in \mathcal{L}$.
  \item A total visit count
        $T_n = \sum_{\ell} c_n(\ell)$.
  \item Depth $d_n$ (distance from the root).
  \item A timestamp $t_n$ recording the last update step.
  \item A transition motor $M_n$ derived from the parent and child
        Context multivectors.
\end{itemize}

Paths from root to leaf encode observed token sequences.  The depth of
the trie grows with the complexity of the data, not with a fixed
parameter budget.

\subsubsection{Lazy Decay}

All counts decay exponentially, but decay is computed lazily: a node's
counts are updated only when it is accessed.  If the current global step
is $s$ and a node was last touched at step $t_n$, the elapsed decay is:
\begin{equation}
  c_n(\ell) \leftarrow c_n(\ell) \cdot \lambda^{s - t_n}
  \label{eq:lazy_decay}
\end{equation}
where $\lambda \in (0,1)$ is the decay factor (default $0.995$).  This
ensures that stale knowledge fades without requiring a global sweep, and
the cost is amortized over accesses.

\subsubsection{Surprisal-Modulated Plasticity}
\label{sec:surprisal_plasticity}

The trie's learning rate is not a fixed hyperparameter.  It is modulated
by \emph{surprisal}: the information-theoretic novelty of each
observation.  For a token $\tau$ observed in context $c$, the surprisal
is:
\begin{equation}
  S(\tau \mid c) = -\log_2 P(\tau \mid c)
  \label{eq:surprisal}
\end{equation}
where $P(\tau \mid c)$ is the trie's current estimate of the next-token
probability given context $c$.  The effective learning rate for this
observation is:
\begin{equation}
  \eta(\tau) = \min\!\bigl(\eta_{\max},\;
    \eta_{\text{base}} + \bar{S} / \sigma_S\bigr)
  \label{eq:adaptive_lr}
\end{equation}
where $\bar{S}$ is the exponential moving average of recent surprisal
values, $\sigma_S$ is its standard deviation, and $\eta_{\text{base}} =
0.1$, $\eta_{\max} = 1.0$ are the floor and ceiling.  Novel inputs
(high surprisal) are learned aggressively; familiar inputs (low
surprisal) are absorbed gently.

\subsubsection{Adaptive Self-Tuning}

Nine parameters that are typically fixed hyperparameters are instead
derived online from signals the trie already computes.  All use
exponential moving averages with smoothing factor $\alpha = 0.05$
(effective window $\approx 20$ observations).  The adaptive mappings
are:

\begin{table}[h]
  \centering
  \caption{Adaptive parameter derivation.  Each parameter is a
           deterministic function of an online signal; no search or
           meta-optimization is required.}
  \label{tab:adaptive_params}
  \begin{tabular}{lll}
    \toprule
    \textbf{Parameter} & \textbf{Signal} & \textbf{Mechanism} \\
    \midrule
    Decay factor $\lambda$
      & Surprisal EMA
      & High surprisal $\to$ lower $\lambda$ (forget faster) \\
    Learning rate $\eta$
      & Per-token surprisal
      & Novel tokens learned more aggressively \\
    Prune threshold
      & Node growth rate
      & Fast growth $\to$ prune harder to control memory \\
    Classification context
      & Classification entropy
      & High entropy $\to$ widen context window \\
    Temperature $\tau$
      & Surprisal EMA
      & Familiar $\to$ explore; novel $\to$ exploit \\
    Beam width
      & Classification confidence
      & Low confidence $\to$ wider search \\
    Interpolation weights
      & Depth hit rate
      & Tracks which $n$-gram depths are predictive \\
    Episodic blend weight
      & Episodic match quality
      & Good episodic hits $\to$ higher blend \\
    Unsupervised threshold
      & Classification accuracy
      & Maturing labels $\to$ require more confidence \\
    \bottomrule
  \end{tabular}
\end{table}

A cold-start guard prevents any adaptive parameter from deviating from
its base value until at least 50 surprisal observations have been
collected, ensuring that the EMA estimates are calibrated before they
influence behavior.

\subsubsection{Classification}

Classification follows a Naive Bayes formulation over interpolated
$n$-gram probabilities.  For an input sequence
$\mathbf{x} = (x_1, \ldots, x_T)$ and label $\ell$, the log-evidence
is:
\begin{equation}
  \log P(\ell \mid \mathbf{x})
    \propto \log P(\ell)
    + \sum_{t=1}^{T} \sum_{d=1}^{D}
      w_d \,\log P_d(x_t \mid x_{t-d:t-1},\, \ell)
  \label{eq:classification}
\end{equation}
where $P_d$ is the $d$-gram probability read from the trie at depth
$d$, $w_d$ is the adaptive interpolation weight for depth $d$, $D$ is
the interpolation suffix depth, and $P(\ell)$ is the label prior from
global class totals.  The system walks the trie depth-first,
accumulating log-evidence per label, and returns the label with the
highest posterior.

\subsubsection{Generation}

Text generation uses beam search with adaptively derived parameters.
At each step, the trie samples candidate next tokens from the
distribution at the current node:
\begin{equation}
  P(\tau \mid c) = \frac{c_n(\tau)}{\sum_{\tau'} c_n(\tau')}
  \label{eq:generation}
\end{equation}
Temperature $\tau$, beam width, and maximum length are not
hyperparameters; they are derived from the trie's adaptive state
(\Cref{tab:adaptive_params}).  Beam candidates are extended while
their cumulative probability exceeds a threshold, producing
completions whose length and diversity adapt to the trie's confidence
in the current context.  Candidate rescoring is two-phase.  First, the
trie-local GF(257) phase vector boosts continuations that interfere
constructively with the current local phase.  Second, the geometric lane
derives a candidate motor from the current Context multivector and the
candidate's deterministic multivector, applies it through the sandwich
product, and boosts candidates whose steered point aligns with the local
attractor:
\begin{equation}
  \text{boost}(\tau)
    = \log\!\left(1 + g \cdot
      \frac{\cos(M_{\tau} X M_{\tau}^{\dagger}, A) +
            \cos(X_{\tau}, A)}{2}\right)
  \label{eq:geometric_continuation_boost}
\end{equation}
where $X$ is the current context coordinate, $X_{\tau}$ is the candidate
coordinate, $A$ is the mean Context multivector over the prediction
path, and $g$ is a fixed gain in the scorer.

\subsubsection{Cross-Domain Procrustes Alignment}

The \texttt{MultimodalCoordinator} maintains independent sensory,
action, and reward tries.  When sensory--action--reward triples are
observed, the coordinator records immutable coactivation statistics:
terminal affinities, terminal IDs, reward residuals, and the sensory,
action, and reward Context multivectors.  Once at least eight paired
sensory/action samples are available, it solves an Orthogonal Procrustes
problem:
\begin{equation}
  R^* = \arg\min_{R^\mathsf{T}R = I} \lVert AR - B \rVert_F
  \label{eq:domain_procrustes}
\end{equation}
where $A$ contains sensory coordinates and $B$ contains action
coordinates.  The resulting rotation is published atomically as an
immutable domain alignment.  During policy inference, sensory
coordinates can be projected into the action manifold before candidate
actions are scored, giving the coordinator a direct geometric path in
addition to affinity and causal-residual evidence.

\subsubsection{Episodic Memory}

A bounded episodic buffer (default capacity 1000) stores recent
token sequences with timestamps and a deterministic geometric coordinate
for each event.  During prediction, the buffer
contributes a $k$-nearest-neighbor tail to the trie's probability
estimates.  The blend weight between trie predictions and episodic
recall is itself adaptive: when episodic matches are high-quality
(low edit distance to the query), the blend weight increases; when
matches are poor, the trie's own estimates dominate.  This provides
fast adaptation to recent context without retraining the trie.

The buffer also exposes a Procrustes re-alignment operation for idle
consolidation.  A caller samples stored events, resolves their current
coordinates in the backing trie, solves a Procrustes rotation from old
episodic coordinates to current trie coordinates, and applies that
rotation to the whole buffer.  This prevents old short-term memories from
drifting away from a manifold that has moved under continued learning.

\subsubsection{Experience-Driven Learning}

The trie does not distinguish training from inference.  When the
\texttt{Predict} method encounters high-surprisal input (average
per-token surprisal above a configurable threshold, default 1.0 bit),
it invokes \texttt{Experience} before classification.  Experience
creates a new concept label if no existing label adequately explains
the input, inserts the sequence into the trie under that label, and
updates all adaptive statistics.  This closed loop means the trie
continuously refines its model as it processes queries.
